\scanpage{178}

\deleted{\uline{\unsure{Prop.}} \unsure{Cor. 3}}

\[
\begin{aligned}
\text{\uuline{Pic}}_{X/S}&=\text{enveloppe exacte de }\text{\uuline{Prépic}}_{X/S}\ (\text{ou }\text{\uuline{Prépic}}'_{X/S})\\
&\text{pour la descente fid plate quasi-compacte}
\end{aligned}
\]

[Il y a un ennui, car un \uline{recouvrement} ouvert d'un préschéma non qu. cpt. n'est en général pas majoré par un morphisme fid plat qu. cpt. Il faut \unsure{donc|sans} dire au lieu de fid plat qu. cpt : \deleted{fid plat après localisation} \uline{morphismes de localisation fidèlement plats}]

Donc on a des morphismes fonctoriels

\diagram{\begin{tikzcd}[ampersand replacement=\&]
\text{\uuline{Pré\deleted{\unsure{P}}pic}}_{X/S}\arrow[r,"\alpha"]\&\text{\uuline{Prépic}}'_{X/S}\arrow[r,"\beta"]\&\text{\uuline{Pic}}_{X/S}
\end{tikzcd}}

Le premier morphisme n'est pas injectif en général [son noyau, si $\underline H^0_{X/S}(\underline O_X)=\underline O_S$, et si cela reste vrai par changmt de base] est le foncteur \uuline{Pic} (absolu) en vertu de prop. 1, (i), \inserted{[qui est nul \unsure{en passant sur un schéma|sur les schémas} \unsure{local|locaux}]} Il est surjectif si on suppose en plus de la condition de [~] que $X/S$ a une section. \deleted{\unsure{Pour avoir la}}

% archive sheet number: 177
